\section{Interpretation and Limitations}
\label{sec:discussion}

The physical bank explains a PMU history through interventions, states and
network hypotheses. Its distinguishing operational choice is to retain a
restored intervention as a possible cause of the continuing response.
This makes recovery compete with a later event instead of treating every
departure from a flat baseline as new. The nominal/mixture comparison
shows a trade-off in development: fewer extras can accompany a late source
decision. Stricter admission eliminates extras but can also reject a real
load. These observations support reporting timely attribution and alarm
burden together.

This construction builds on established trajectory-based diagnosis,
multiple-model filtering and sequential inference. Warner's physical
trajectory scan already handles unknown onset and sparse remote PMUs;
Adaptive Matrix CuSum provides error/delay theory for its statistical
setting. We do not inherit those guarantees or claim numerical superiority
over methods not executed on common data. The present admission ablation
also does not isolate the causal benefit of every recovery-state operation.

The executed scope is a known catalog of nine loads and nine lines, with
an intervention-free startup. Missing PMUs, unknown event families,
source-excluded transfer and field operation are not tested. F01's
measurement map includes the provisional BUS5 current terminal and excludes
BUS39 current angle. Agreement with this research model does not establish
the original producer's CT placement or exact replication of RAW0001.
Local predictors, pruning, shared covariance approximations and plug-in
startup estimates limit the interpretation of the masses as probabilities.

The ExtraTrees comparison concerns two specific fixed causal adapters.
Their load-transfer failure is not evidence that supervised learning is
incapable of recognizing loads, and its score cannot be substituted for a
historical multi-family benchmark. Chen's F01 adapter has DC/white-injection
assumptions that do not capture all instrumental effects; its failures
must be reported with that mismatch.

Physics-derived stochastic priors offer a coherent direction for future
prediction uncertainty \cite{ensaf2023kernel,tartakovsky2023gp}, while
joint DAE parameter estimation addresses another task \cite{albustami2026dae}.
Such extensions would require separate validation. They are not components
silently added to the method evaluated here. The current scientific claim
rests on the specified recovery-aware mechanism and the evidence produced
by its fixed experimental contract.

A specific successor question is whether context should change the event
signature itself. For pre-event context $\eta$, a mixed sensitivity
$E_i=\partial D_i/\partial\eta$ changes the predicted effect of source $i$,
rather than merely shifting the nominal trajectory. Under a local Gaussian
context approximation, its predictive covariance would be
\begin{equation}
 V_h(a)=\Sigma_W+J_{\eta,h}(a)P_{\eta,\tau}J_{\eta,h}(a)^\top.
 \label{eq:context-conditioned-extension}
\end{equation}
The determinant of $V_h(a)$ must remain in the evidence. Testing this
extension requires matched recentering-only and mixed-sensitivity controls,
calibrated pre-event uncertainty and new physical trajectories. The present
R11 results do not validate that joint multiple-source algorithm.

A decisive test would hold observations, priors and horizon fixed when
comparing nominal, recentered and context-conditioned signatures. Varying
horizon separately would identify gains from additional data.
Nonlinear-response controls would distinguish surrogate error from
confusion among physical alternatives. These successor experiments must
measure source decisions and uncertainty calibration; algebraic
equivalence or lower prediction error alone cannot establish either benefit.
